\documentclass[aspectratio=169]{beamer}

% Use Sorbonne theme
\usepackage{sorbonne_theme}

% Title page information
\title{Econometrics}
\subtitle{Ridge and Lasso Regression}
\author{Jose Angel Garcia Sanchez}
\institute{Université Paris 1 Panthéon-Sorbonne}
\date{} % left empty on purpose: beamer would otherwise print today's date
% Set the logo
\sorbonnelogo{imagen.png}

\begin{document}

% Title page (without logo in corner)
{
\setbeamertemplate{logo}{}
\begin{frame}
\titlepage
\end{frame}
}

% Table of contents
\begin{frame}{Outline}
\tableofcontents
\end{frame}

% ===============================
% SECTION 1: Introduction to Regularization
% ===============================
\section{Introduction to Regularization}

\begin{frame}{Motivation}
\begin{itemize}
    \item OLS minimizes the sum of squared residuals: $\min_{\beta} \sum_{i=1}^{n}(y_i - x_i'\beta)^2$
    \item OLS works well when $n \gg p$ (many observations, few variables)
    \item But OLS can fail when:
    \begin{itemize}
        \item $p$ is large relative to $n$ (high-dimensional data); if $p > n$, $X'X$ is not even invertible
        \item There is multicollinearity among regressors
        \item We want to improve out-of-sample prediction
    \end{itemize}
    \item Solution: add a \textit{penalty term} to the objective function
    \item This is called \textit{regularization} or \textit{shrinkage}
\end{itemize}
\end{frame}

\begin{frame}{The Bias-Variance Trade-off}
\begin{itemize}
    \item Recall that the Mean Squared Error of an estimator can be decomposed: $MSE = Bias^2 + Variance$
    \item OLS is unbiased but can have high variance, especially when:
    \begin{itemize}
        \item Predictors are highly correlated (multicollinearity)
        \item The number of predictors is large
    \end{itemize}
    \item Regularization introduces some bias but reduces variance
    \item If the reduction in variance exceeds the increase in squared bias, we get a lower MSE
    \item This improves prediction accuracy, especially out-of-sample
\end{itemize}
\end{frame}

\begin{frame}{The General Penalized Regression Framework}
\begin{itemize}
    \item General form: $\min_{\beta} \left\{ \sum_{i=1}^{n}(y_i - x_i'\beta)^2 + \lambda \cdot P(\beta) \right\}$
    \item $\lambda \geq 0$ is the \textit{tuning parameter} (or regularization parameter)
    \item $P(\beta)$ is the \textit{penalty function}
    \item When $\lambda = 0$: we get OLS
    \item When $\lambda \rightarrow \infty$: all coefficients shrink to zero
    \item Different choices of $P(\beta)$ give different estimators:
    \begin{itemize}
        \item Ridge: $P(\beta) = \sum_{j=1}^{p}\beta_j^2$ (L2 penalty)
        \item Lasso: $P(\beta) = \sum_{j=1}^{p}|\beta_j|$ (L1 penalty)
        \item Elastic Net: combination of both
    \end{itemize}
\end{itemize}
\end{frame}

\begin{frame}{Important Preliminary: Standardization}
\begin{itemize}
    \item Regularization penalizes the size of coefficients
    \item But coefficient size depends on the scale of variables
    \item A variable measured in millions will have a tiny coefficient compared to one measured in units
    \item Solution: \textit{standardize} all variables before estimation
    \item For each variable $x_j$: $\tilde{x}_j = \frac{x_j - \bar{x}_j}{s_{x_j}}$
    \item After standardization, all variables have mean 0 and standard deviation 1
    \item The intercept is typically not penalized
    \item After estimation, coefficients can be transformed back to the original scale if needed
\end{itemize}
\end{frame}

% ===============================
% SECTION 2: Ridge Regression
% ===============================
\section{Ridge Regression}

\begin{frame}{Ridge Regression: Definition}
\begin{itemize}
    \item Ridge regression adds an L2 penalty to the OLS objective:
    \[\hat{\beta}^{Ridge} = \arg\min_{\beta} \left\{ \sum_{i=1}^{n}(y_i - x_i'\beta)^2 + \lambda \sum_{j=1}^{p}\beta_j^2 \right\}\]
    \item Equivalently, in matrix notation:
    \[\hat{\beta}^{Ridge} = \arg\min_{\beta} \left\{ (Y - X\beta)'(Y - X\beta) + \lambda \beta'\beta \right\}\]
    \item The penalty term $\lambda \sum_{j=1}^{p}\beta_j^2$ shrinks coefficients toward zero
    \item Large coefficients are penalized more heavily (quadratic penalty)
\end{itemize}
\end{frame}

\begin{frame}{Ridge Regression: Closed-Form Solution}
\begin{itemize}
    \item Taking the first-order condition:
    \[\frac{\partial}{\partial \beta}\left[(Y - X\beta)'(Y - X\beta) + \lambda \beta'\beta\right] = 0\]
    \item We get: $-2X'Y + 2X'X\beta + 2\lambda\beta = 0$
    \item Solving for $\beta$:
    \[\hat{\beta}^{Ridge} = (X'X + \lambda I)^{-1}X'Y\]
    \item Compare with OLS: $\hat{\beta}^{OLS} = (X'X)^{-1}X'Y$
    \item The term $\lambda I$ is added to $X'X$ before inversion
    \item This ensures $(X'X + \lambda I)$ is always invertible for $\lambda > 0$, even when $X'X$ is singular
\end{itemize}
\end{frame}

\begin{frame}{Ridge Regression: Properties (1)}
\begin{itemize}
    \item The Ridge estimator is biased: $E(\hat{\beta}^{Ridge}) \neq \beta$
    \item The bias is: $Bias(\hat{\beta}^{Ridge}) = -\lambda(X'X + \lambda I)^{-1}\beta$
    \item The variance is: $Var(\hat{\beta}^{Ridge}) = \sigma^2(X'X + \lambda I)^{-1}X'X(X'X + \lambda I)^{-1}$
    \item As $\lambda$ increases:
    \begin{itemize}
        \item Bias increases (coefficients shrink more toward zero)
        \item Variance decreases (estimates become more stable)
    \end{itemize}
    \item There exists an optimal $\lambda^*$ that minimizes MSE, and it is always possible to beat OLS: some $\lambda > 0$ gives a lower MSE than OLS (Hoerl and Kennard, 1970)
\end{itemize}
\end{frame}

\begin{frame}{Ridge Regression: Properties (2)}
\begin{itemize}
    \item Ridge does NOT set coefficients exactly to zero
    \item All predictors remain in the model (no variable selection)
    \item Ridge is particularly useful when:
    \begin{itemize}
        \item Predictors are highly correlated (multicollinearity)
        \item We believe all predictors are relevant
        \item We want to stabilize estimates
    \end{itemize}
    \item The ridge penalty can be viewed as a Bayesian prior: $\beta \sim N(0, (\sigma^2/\lambda) I)$
\end{itemize}
\end{frame}

\begin{frame}{Ridge Regression: Geometric Interpretation}
\begin{itemize}
    \item Ridge regression can be written as a constrained optimization:
    \[\min_{\beta} \sum_{i=1}^{n}(y_i - x_i'\beta)^2 \quad \text{subject to} \quad \sum_{j=1}^{p}\beta_j^2 \leq t\]
    \item The constraint region is a ball (a disk in 2 dimensions)
    \item The OLS solution is at the center of the elliptical contours of the RSS
    \item The ridge solution is where the elliptical contours first touch the ball
    \item Because the ball has no corners, the solution typically has all coefficients non-zero
    \item Smaller $t$ (equivalently, larger $\lambda$) means more shrinkage
\end{itemize}
\end{frame}

\begin{frame}{Ridge and Multicollinearity}
\begin{itemize}
    \item Recall: with perfect multicollinearity, $X'X$ is singular and OLS fails
    \item With near-perfect multicollinearity, $(X'X)^{-1}$ has very large elements
    \item This leads to inflated variances and unstable OLS estimates
    \item Ridge regression solves this by adding $\lambda I$ to $X'X$
    \item Even if $X'X$ is singular, $(X'X + \lambda I)$ is always invertible for $\lambda > 0$
    \item The eigenvalues of $(X'X + \lambda I)$ are $d_j + \lambda$, where $d_j$ are the eigenvalues of $X'X$
    \item Small eigenvalues (source of instability) become at least $\lambda$ instead of nearly zero
\end{itemize}
\end{frame}

% ===============================
% SECTION 3: Lasso Regression
% ===============================
\section{Lasso Regression}

\begin{frame}{Lasso Regression: Definition}
\begin{itemize}
    \item Lasso = Least Absolute Shrinkage and Selection Operator (Tibshirani, 1996)
    \item Lasso adds an L1 penalty to the OLS objective:
    \[\hat{\beta}^{Lasso} = \arg\min_{\beta} \left\{ \sum_{i=1}^{n}(y_i - x_i'\beta)^2 + \lambda \sum_{j=1}^{p}|\beta_j| \right\}\]
    \item The key difference from Ridge: absolute values instead of squares
    \item This seemingly small change has profound implications
    \item Lasso performs both shrinkage AND variable selection
\end{itemize}
\end{frame}

\begin{frame}{Lasso: No Closed-Form Solution}
\begin{itemize}
    \item Unlike Ridge, Lasso has no closed-form solution
    \item The absolute value function $|\beta_j|$ is not differentiable at $\beta_j = 0$
    \item Optimization requires iterative algorithms:
    \begin{itemize}
        \item Coordinate descent (most common)
        \item LARS (Least Angle Regression)
        \item Proximal gradient methods
    \end{itemize}
    \item With an orthonormal design ($X'X = I$), each coefficient is a \textit{soft-thresholding} of the OLS one:
    \[\hat{\beta}_j^{Lasso} = \text{sign}(\hat{\beta}_j^{OLS})\left(|\hat{\beta}_j^{OLS}| - \lambda/2\right)_+\]
    \item Where $(x)_+ = \max(0, x)$: small OLS coefficients are set exactly to zero
    \item Compare Ridge in the same case: $\hat{\beta}_j^{Ridge} = \hat{\beta}_j^{OLS}/(1 + \lambda)$, a proportional shrinkage that never reaches zero
\end{itemize}
\end{frame}

\begin{frame}{Lasso: Variable Selection Property}
\begin{itemize}
    \item The key feature of Lasso: it sets some coefficients exactly to zero
    \item This performs automatic variable selection
    \item As $\lambda$ increases, more coefficients become exactly zero
    \item The model becomes more sparse (fewer non-zero coefficients)
    \item This is very useful when:
    \begin{itemize}
        \item We have many potential predictors
        \item We believe only a subset are truly relevant
        \item We want an interpretable, parsimonious model
    \end{itemize}
    \item Lasso identifies which variables matter and estimates their effects simultaneously
\end{itemize}
\end{frame}

\begin{frame}{Lasso: Geometric Interpretation}
\begin{itemize}
    \item Lasso can be written as a constrained optimization:
    \[\min_{\beta} \sum_{i=1}^{n}(y_i - x_i'\beta)^2 \quad \text{subject to} \quad \sum_{j=1}^{p}|\beta_j| \leq t\]
    \item The constraint region is a diamond (or cross-polytope in higher dimensions)
    \item The diamond has corners on the axes
    \item The elliptical RSS contours often touch the diamond at a corner
    \item At a corner, one or more coefficients are exactly zero
    \item This is why Lasso produces sparse solutions
\end{itemize}
\end{frame}

\begin{frame}{Lasso: Properties}
\begin{itemize}
    \item The Lasso estimator is biased (like Ridge)
    \item Lasso is consistent for variable selection under certain conditions
    \item However, Lasso has some limitations:
    \begin{itemize}
        \item With highly correlated predictors, Lasso tends to select one and ignore the others
        \item When $p > n$, Lasso selects at most $n$ variables
        \item Lasso estimates for selected variables are biased toward zero
    \end{itemize}
    \item These limitations motivate extensions like Elastic Net and Adaptive Lasso
\end{itemize}
\end{frame}

\begin{frame}{Comparing Ridge and Lasso}
\begin{center}
\begin{tabular}{l|c|c}
\hline
Property & Ridge & Lasso \\
\hline
Penalty & $\sum \beta_j^2$ (L2) & $\sum |\beta_j|$ (L1) \\
Closed-form solution & Yes & No \\
Variable selection & No & Yes \\
Handles multicollinearity & Yes & Partially \\
Correlated predictors & Shrinks together & Selects one \\
When $p > n$ & Works & Selects $\leq n$ \\
Bayesian interpretation & Normal prior & Laplace prior \\
\hline
\end{tabular}
\end{center}
\end{frame}

% ===============================
% SECTION 4: Elastic Net
% ===============================
\section{Elastic Net}

\begin{frame}{Elastic Net: Combining Ridge and Lasso}
\begin{itemize}
    \item Elastic Net (Zou and Hastie, 2005) combines L1 and L2 penalties:
    \[\hat{\beta}^{EN} = \arg\min_{\beta} \left\{ \sum_{i=1}^{n}(y_i - x_i'\beta)^2 + \lambda_1 \sum_{j=1}^{p}|\beta_j| + \lambda_2 \sum_{j=1}^{p}\beta_j^2 \right\}\]
    \item Alternative parameterization with mixing parameter $\alpha \in [0,1]$:
    \[\hat{\beta}^{EN} = \arg\min_{\beta} \left\{ \sum_{i=1}^{n}(y_i - x_i'\beta)^2 + \lambda \left[ \alpha \sum_{j=1}^{p}|\beta_j| + (1-\alpha) \sum_{j=1}^{p}\beta_j^2 \right] \right\}\]
    \item $\alpha = 1$: Lasso; $\alpha = 0$: Ridge
\end{itemize}
\end{frame}

\begin{frame}{Elastic Net: Properties}
\begin{itemize}
    \item Combines the best of both worlds:
    \begin{itemize}
        \item Variable selection (from Lasso)
        \item Grouping effect (from Ridge)
    \end{itemize}
    \item With correlated predictors, Elastic Net tends to select groups of correlated variables together
    \item Can select more than $n$ variables when $p > n$
    \item Particularly useful when:
    \begin{itemize}
        \item Predictors are highly correlated
        \item We want both sparsity and stability
        \item The number of predictors exceeds the sample size
    \end{itemize}
\end{itemize}
\end{frame}

% ===============================
% SECTION 5: Choosing the Tuning Parameter
% ===============================
\section{Choosing the Tuning Parameter}

\begin{frame}{The Role of $\lambda$}
\begin{itemize}
    \item The tuning parameter $\lambda$ controls the amount of regularization
    \item $\lambda = 0$: no penalty, we get OLS
    \item $\lambda \rightarrow \infty$: maximum penalty, all coefficients go to zero
    \item Intermediate $\lambda$: trade-off between bias and variance
    \item How to choose $\lambda$? Several approaches:
    \begin{itemize}
        \item Cross-validation (most common)
        \item Information criteria (AIC, BIC)
        \item Theoretical considerations
    \end{itemize}
\end{itemize}
\end{frame}

\begin{frame}{Cross-Validation (1)}
\begin{itemize}
    \item K-fold cross-validation is the standard approach
    \item Procedure:
    \begin{enumerate}
        \item Divide data into $K$ roughly equal parts (folds)
        \item For each fold $k = 1, \ldots, K$:
        \begin{itemize}
            \item Fit the model on all data except fold $k$
            \item Compute the prediction error on fold $k$
        \end{itemize}
        \item Average the $K$ prediction errors
    \end{enumerate}
    \item Repeat for a grid of $\lambda$ values
    \item Choose the $\lambda$ that minimizes the cross-validation error
    \item Common choices: $K = 5$ or $K = 10$
\end{itemize}
\end{frame}

\begin{frame}{Cross-Validation (2)}
\begin{itemize}
    \item The CV error as a function of $\lambda$ typically has a U-shape
    \item Small $\lambda$: high variance, low bias (overfitting)
    \item Large $\lambda$: low variance, high bias (underfitting)
    \item Two common choices:
    \begin{itemize}
        \item $\lambda_{min}$: minimizes the CV error
        \item $\lambda_{1se}$: largest $\lambda$ whose CV error is within 1 standard error of the minimum
    \end{itemize}
    \item $\lambda_{1se}$ gives a more parsimonious model (more regularization)
    \item This is the ``one-standard-error rule''
\end{itemize}
\end{frame}

\begin{frame}{Information Criteria}
\begin{itemize}
    \item Alternative to cross-validation: use information criteria
    \item For Lasso, we can use:
    \begin{itemize}
        \item AIC: $AIC = n \log(RSS/n) + 2\,df$
        \item BIC: $BIC = n \log(RSS/n) + \log(n) \cdot df$
    \end{itemize}
    \item The degrees of freedom of the Lasso can be estimated by the number of non-zero coefficients
    \item For Ridge: $df = \sum_{j=1}^{p} \frac{d_j}{d_j + \lambda}$, where $d_j$ are the eigenvalues of $X'X$
    \item BIC tends to select sparser models than AIC
    \item Cross-validation is generally preferred in practice
\end{itemize}
\end{frame}

% ===============================
% SECTION 6: Inference and Extensions
% ===============================
\section{Inference and Extensions}

\begin{frame}{Inference After Selection}
\begin{itemize}
    \item Standard inference (t-tests, confidence intervals) is invalid after Lasso selection
    \item Why? The selection process introduces additional randomness
    \item Naive standard errors are too small (underestimate uncertainty)
    \item Confidence intervals have incorrect coverage
    \item Several approaches have been developed:
    \begin{itemize}
        \item Sample splitting
        \item Selective inference
        \item Debiased/Desparsified Lasso
        \item Bootstrap methods
        \item Post-double-selection (Belloni, Chernozhukov and Hansen, 2014), for the effect of one variable of interest in the presence of many controls
    \end{itemize}
\end{itemize}
\end{frame}

\begin{frame}{Sample Splitting}
\begin{itemize}
    \item Simple approach to valid inference:
    \begin{enumerate}
        \item Split the sample randomly into two parts
        \item Use the first part for variable selection (Lasso)
        \item Use the second part for estimation and inference (OLS on selected variables)
    \end{enumerate}
    \item Advantages: simple, valid inference
    \item Disadvantages:
    \begin{itemize}
        \item Loses statistical power (uses only half the data for each step)
        \item Results depend on the random split
    \end{itemize}
    \item Can average over multiple random splits
\end{itemize}
\end{frame}

\begin{frame}{Post-Lasso OLS}
\begin{itemize}
    \item Also called ``Post-Selection OLS'' or ``OLS post Lasso''
    \item Procedure:
    \begin{enumerate}
        \item Run Lasso to select variables (non-zero coefficients)
        \item Run OLS using only the selected variables
    \end{enumerate}
    \item The OLS estimates are less biased than Lasso estimates
    \item Removes the shrinkage bias for selected variables
    \item Caution: standard errors from this OLS are still not valid for inference
    \item Valid inference requires additional corrections
\end{itemize}
\end{frame}

\begin{frame}{Adaptive Lasso}
\begin{itemize}
    \item Standard Lasso penalizes all coefficients equally
    \item This can lead to too much bias for large coefficients
    \item Adaptive Lasso (Zou, 2006) uses weighted penalties:
    \[\hat{\beta}^{ALasso} = \arg\min_{\beta} \left\{ \sum_{i=1}^{n}(y_i - x_i'\beta)^2 + \lambda \sum_{j=1}^{p}w_j|\beta_j| \right\}\]
    \item Weights: $w_j = 1/|\hat{\beta}_j^{init}|^\gamma$ for some initial estimator $\hat{\beta}^{init}$ (e.g.\ OLS or Ridge)
    \item Typically $\gamma = 1$ or $\gamma = 2$
    \item Variables with small initial estimates are penalized more
    \item Under some conditions, Adaptive Lasso has the ``oracle property'': asymptotically, it selects the correct model and estimates its coefficients as if the true model were known
\end{itemize}
\end{frame}

% ===============================
% SECTION 7: Practical Implementation
% ===============================
\section{Practical Implementation}

\begin{frame}{Implementation in Software}
\begin{itemize}
    \item Ridge and Lasso are implemented in all major statistical software:
    \begin{itemize}
        \item R: \texttt{glmnet} package (gold standard)
        \item Python: \texttt{scikit-learn} (\texttt{Ridge}, \texttt{Lasso}, \texttt{ElasticNet}, and \texttt{RidgeCV}, \texttt{LassoCV}, \texttt{ElasticNetCV} with built-in cross-validation)
        \item Stata: \texttt{lasso}, \texttt{elasticnet} commands
    \end{itemize}
    \item \texttt{glmnet} uses coordinate descent, computes the entire solution path (all values of $\lambda$) efficiently, and standardizes the variables by default
    \item scikit-learn does \textbf{not} standardize: use \texttt{make\_pipeline(StandardScaler(), LassoCV())}
    \item scikit-learn calls $\lambda$ \texttt{alpha}, and divides the Lasso and Elastic Net RSS by $2n$: its \texttt{alpha} is $\lambda/(2n)$ in our notation (but $\lambda$ for \texttt{Ridge}); our mixing parameter $\alpha$ is its \texttt{l1\_ratio}
\end{itemize}
\end{frame}

\begin{frame}{Practical Recommendations}
\begin{itemize}
    \item Always standardize predictors before regularization
    \item Use cross-validation to choose $\lambda$
    \item Consider the one-standard-error rule for more parsimony
    \item If predictors are highly correlated, consider Elastic Net
    \item For inference, use appropriate methods (not naive standard errors)
    \item Check model performance on held-out test data
    \item Compare with OLS on a subset of variables (sanity check)
    \item Remember: regularization is primarily for prediction, not causal inference (unless using methods designed for it, such as post-double-selection)
\end{itemize}
\end{frame}

\begin{frame}{When to Use What?}
\begin{itemize}
    \item \textbf{OLS}: $n \gg p$, no multicollinearity, focus on inference
    \item \textbf{Ridge}: multicollinearity, all predictors believed relevant, prediction focus
    \item \textbf{Lasso}: many predictors, only some believed relevant, want variable selection
    \item \textbf{Elastic Net}: many correlated predictors, want both selection and grouping
    \item \textbf{Adaptive Lasso}: want consistent variable selection with oracle properties
    \item In practice: try several methods and compare via cross-validation
\end{itemize}
\end{frame}

\begin{frame}{Summary}
\begin{itemize}
    \item Regularization adds a penalty to the OLS objective function
    \item Ridge (L2): shrinks coefficients, handles multicollinearity, no selection
    \item Lasso (L1): shrinks and selects variables, produces sparse models
    \item Elastic Net: combines Ridge and Lasso benefits
    \item The tuning parameter $\lambda$ is chosen via cross-validation
    \item Standard inference is invalid after Lasso selection
    \item These methods are powerful tools for prediction in high-dimensional settings
    \item But remember: regularization introduces bias, which may not be desirable for causal inference
\end{itemize}
\end{frame}

\end{document}
